\section{Discussion and Limitations}
\label{sec:discussion}

\paragraph{What worked.} Two changes give most of the gain over the public baseline: using
the anatomical correspondence that the mesh already provides, and estimating the symmetry
axis instead of assuming it. The mirror test is cheap, needs no labels, and gives the
system a landmark-reliability signal that MediaPipe does not expose. The gate turns many
silent failures into explicit refusals.

\paragraph{What did not.} The mirror test cannot detect an occluded eye that the detector
fills in consistently. On clean photos the gate does not predict small asymmetry errors
better than chance, because those errors are dominated by symmetric landmark noise and
by the detector's shrinkage towards symmetry.

\paragraph{Limitations.}
\begin{itemize}[leftmargin=*,itemsep=1pt,topsep=2pt]
\item \textbf{2D only.} A photo is a projection. Yaw and pitch change the measured value
(E8), and the gate can only limit, not remove, that effect. Camera perspective at selfie
distance is not modelled.
\item \textbf{Reference, not ground truth.} $A_{ref}$ uses expert landmarks, which contain
annotator noise, and is subject to the same projection effects. We have no human
symmetry ratings, and a validated perceptual or clinical scale was out of scope.
\item \textbf{Dataset.} LFPW is small (224 test faces), in the wild and mostly non-frontal,
and each identity appears once, so the pose strata are not paired. MediaPipe's training
data is not public and could overlap with LFPW, so its NME here is not an independent
benchmark claim. We use it only to compare mappings and gate signals.
\item \textbf{Thresholds} are percentiles of one training distribution. A different camera
or population would need re-calibration (\texttt{calibrate\_gate.py}).
\item \textbf{Expression} is not controlled. A smile is asymmetric by nature, so the app
asks for a neutral expression.
\end{itemize}

\paragraph{Ethics and privacy.}
\label{sec:ethics}
The output is a geometric indicator. It says nothing about health, identity or
attractiveness, and the app states this. Dataset images are used under their research
licence and are neither redistributed in the repository nor reproduced in this report.
Figures draw landmark geometry on a blank canvas. The app processes uploads in memory and
stores nothing.

\section{Conclusion and Future Work}
We built and evaluated a quality-aware facial asymmetry system that reports a measurement
only when the image supports it. On LFPW it tracks an annotation-based reference much more
closely than a re-implemented public baseline ($\rho$ \rhoMthree{} vs.\ \rhoMzero{}). Its
mirror test flags unreliable landmarks without labels, and its gate cuts silent failures
under degradation from \silentPoolOne\% to \silentPoolThree\%. The study also gives a
practical warning: in uncontrolled photos, head yaw explains most of the measured 2D
asymmetry. Future work: frontalisation with the detector's 3D landmarks, evaluated on a
dataset with ground-truth pose (\eg AFLW2000-3D~\cite{zhu2016aflw2000}); appearance-based
occlusion detection; WFLW's attribute subsets; and a small consented, multi-pose capture of
the same people with human symmetry ratings.
